% GENERATED FILE. Edit canonical lessons, then run npm run generate:books.
% canonical-source-hash: fnv1a64:4936291c31fca20f
% canonical-lessons: ES-C771-sabor, ES-C771-picante, ES-C771-salado, ES-C771-crudo, ES-C771-oler

\chapter{Por el sabor}
\label{ch:es-por-el-sabor}
\hlchaptermodality{771}

\begin{chapteropening}
\textbf{By the end of this chapter:} \emph{I can say why I like or dislike a dish: its taste, whether it is hot, salty or raw, and how it smells.}

Complete the last lesson of chapter 771: I can say why I like or dislike a dish: its taste, whether it is hot, salty or raw, and how it smells.
\end{chapteropening}

\section[el sabor]{el sabor — the taste, the flavour}

\label{lesson:ES-C771-sabor}



\begin{quote}
Before the new one: say the Spanish for the butcher's, then the Spanish for the greengrocer's.
\end{quote}

\subsection*{You'll want to know: el sabor}
\textbf{el sabor} — \textquotedblleft{}the taste, the flavour\textquotedblright{}.

\begin{quote}
Me gusta por el sabor, no por el precio.
\end{quote}

\begin{quote}
\emph{I like it for the taste, not the price.}
\end{quote}

\emph{¿De qué sabor?} — which flavour? \emph{Saber a} is to taste of: \emph{Sabe a limón.}

\begin{cousinweb}
\emph{sabor} is inherited from Latin \emph{sapōrem}, from \emph{sapor} 'a taste, a flavour' (EtymDB), built on \emph{sapere} 'to taste', the verb behind \emph{saber}.
\end{cousinweb}

\begin{grammarlens}[title={What you've built}]
One more word for telling, explaining and describing at B1.
\end{grammarlens}

\subsection*{Your turn}
\begin{itemize}
\raggedright
  \item \emph{Say it:} \emph{el sabor}
  \item \emph{Say it:} \emph{el sabor}, once more
  \item \emph{Say it:} a sentence of your own with \emph{el sabor}
  \item \emph{Recall:} say the Spanish for the butcher's, then the Spanish for the greengrocer's, then say \emph{el sabor} again
\end{itemize}

\begin{tcolorbox}[breakable,colback=teal!4,colframe=teal!35!black,title={Before you move on}]
What does el sabor mean? (\textquotedblleft{}The taste, the flavour\textquotedblright{}.) Say it once more.
\end{tcolorbox}
\section[picante]{picante — spicy, hot}

\label{lesson:ES-C771-picante}



\begin{quote}
Before the new one: say the Spanish for the greengrocer's, then the Spanish for the taste.
\end{quote}

\subsection*{You'll want to know: picante}
\textbf{picante} — \textquotedblleft{}spicy, hot\textquotedblright{}.

\begin{quote}
No lo como porque es demasiado picante.
\end{quote}

\begin{quote}
\emph{I don't eat it because it is too spicy.}
\end{quote}

\emph{Picante} is spicy-hot; hot in temperature is \emph{caliente}.

\begin{cousinweb}
\emph{picante} is built on \emph{picar} 'to sting, to prick'. \emph{Picar} is usually traced to spoken Latin \emph{*pīccāre} 'to sting, to strike', the form EtymDB gives for its French cousin \emph{piquer}. Spicy food stings the tongue.
\end{cousinweb}

\begin{grammarlens}[title={What you've built}]
One more word for telling, explaining and describing at B1.
\end{grammarlens}

\subsection*{Your turn}
\begin{itemize}
\raggedright
  \item \emph{Say it:} \emph{picante}
  \item \emph{Say it:} \emph{picante}, once more
  \item \emph{Say it:} a sentence of your own with \emph{picante}
  \item \emph{Recall:} say the Spanish for the greengrocer's, then the Spanish for the taste, then say \emph{picante} again
\end{itemize}

\begin{tcolorbox}[breakable,colback=teal!4,colframe=teal!35!black,title={Before you move on}]
What does picante mean? (\textquotedblleft{}Spicy, hot\textquotedblright{}.) Say it once more.
\end{tcolorbox}
\section[salado]{salado — salty; savoury}

\label{lesson:ES-C771-salado}



\begin{quote}
Before the new one: say the Spanish for the taste, then the Spanish for spicy.
\end{quote}

\subsection*{You'll want to know: salado}
\textbf{salado} — \textquotedblleft{}salty; savoury\textquotedblright{}.

\begin{quote}
No me gusta la sopa porque está muy salada.
\end{quote}

\begin{quote}
\emph{I don't like the soup because it is very salty.}
\end{quote}

In Spain \emph{salado} can also mean witty or charming; in parts of Latin America, unlucky.

\begin{cousinweb}
\emph{salado} is built on the participle of \emph{salar} 'to salt', from \emph{sal}, which is inherited from Latin \emph{sāl} 'salt' (Lewis \& Short).
\end{cousinweb}

\begin{grammarlens}[title={What you've built}]
One more word for telling, explaining and describing at B1.
\end{grammarlens}

\subsection*{Your turn}
\begin{itemize}
\raggedright
  \item \emph{Say it:} \emph{salado}
  \item \emph{Say it:} \emph{salado}, once more
  \item \emph{Say it:} a sentence of your own with \emph{salado}
  \item \emph{Recall:} say the Spanish for the taste, then the Spanish for spicy, then say \emph{salado} again
\end{itemize}

\begin{tcolorbox}[breakable,colback=teal!4,colframe=teal!35!black,title={Before you move on}]
What does salado mean? (\textquotedblleft{}Salty; savoury\textquotedblright{}.) Say it once more.
\end{tcolorbox}
\section[crudo]{crudo — raw; uncooked}

\label{lesson:ES-C771-crudo}



\begin{quote}
Before the new one: say the Spanish for spicy, then the Spanish for salty.
\end{quote}

\subsection*{You'll want to know: crudo}
\textbf{crudo} — \textquotedblleft{}raw; uncooked\textquotedblright{}.

\begin{quote}
No como sushi porque no me gusta el pescado crudo.
\end{quote}

\begin{quote}
\emph{I don't eat sushi because I don't like raw fish.}
\end{quote}

\emph{El petróleo crudo} — crude oil. \emph{Lo tiene crudo} (Spain, informal) — he's got it tough.

\begin{cousinweb}
\emph{crudo} is inherited from Latin \emph{crūdus} 'raw, uncooked, bloody' (EtymDB), through old Spanish \emph{crudo}.
\end{cousinweb}

\begin{grammarlens}[title={What you've built}]
One more word for telling, explaining and describing at B1.
\end{grammarlens}

\subsection*{Your turn}
\begin{itemize}
\raggedright
  \item \emph{Say it:} \emph{crudo}
  \item \emph{Say it:} \emph{crudo}, once more
  \item \emph{Say it:} a sentence of your own with \emph{crudo}
  \item \emph{Recall:} say the Spanish for spicy, then the Spanish for salty, then say \emph{crudo} again
\end{itemize}

\begin{tcolorbox}[breakable,colback=teal!4,colframe=teal!35!black,title={Before you move on}]
What does crudo mean? (\textquotedblleft{}Raw; uncooked\textquotedblright{}.) Say it once more.
\end{tcolorbox}
\section[oler]{oler — to smell}

\label{lesson:ES-C771-oler}



\begin{quote}
Before the new one: say the Spanish for salty, then the Spanish for raw.
\end{quote}

\subsection*{You'll want to know: oler}
\textbf{oler} — \textquotedblleft{}to smell\textquotedblright{}.

\begin{quote}
No la como porque huele mal.
\end{quote}

\begin{quote}
\emph{I don't eat it because it smells bad.}
\end{quote}

An \emph{h} is added where it changes: \emph{huelo}, \emph{hueles}. \emph{Oler a} — to smell of.

\begin{cousinweb}
\emph{oler} is inherited from Latin \emph{olēre} 'to give off a smell' (EtymDB). The stressed \emph{o} became \emph{ue}, and Spanish writes an \emph{h} in front: \emph{huele}.
\end{cousinweb}

\begin{grammarlens}[title={What you've built}]
One more word for telling, explaining and describing at B1.
\end{grammarlens}

\subsection*{Your turn}
\begin{itemize}
\raggedright
  \item \emph{Say it:} \emph{oler}
  \item \emph{Say it:} \emph{oler}, once more
  \item \emph{Say it:} a sentence of your own with \emph{oler}
  \item \emph{Recall:} say the Spanish for salty, then the Spanish for raw, then say \emph{oler} again
\end{itemize}

\begin{tcolorbox}[breakable,colback=teal!4,colframe=teal!35!black,title={Before you move on}]
What does oler mean? (\textquotedblleft{}To smell\textquotedblright{}.) Say it once more.
\end{tcolorbox}
